\documentclass[11pt]{article}
\usepackage[utf8]{inputenc}
\usepackage[T1]{fontenc}
\usepackage{amsmath}
\usepackage{booktabs}
\usepackage{graphicx}
\usepackage{microtype}
\usepackage{geometry}
\usepackage{hyperref}
\usepackage{xcolor}
\usepackage{enumitem}
\geometry{margin=1in}
\hypersetup{colorlinks=true,linkcolor=blue,citecolor=blue,urlcolor=blue}

\title{SignLearn: Transfer-Initialized Motion Fusion and Selective Prediction\\
for a 50-Class Isolated ASL Practice Prototype}
\author{Anonymous working draft}
\date{Draft manuscript --- August 2026}

\begin{document}
\maketitle

\begin{abstract}
This paper presents SignLearn, a reproducible research prototype for classifying one of 50
isolated American Sign Language (ASL) labels from a short webcam or uploaded video. The system
combines MediaPipe Holistic landmarks with a 492-dimensional per-frame representation containing
normalized appearance, first- and second-order motion, and compact body--hand--face geometry.
A bidirectional long short-term memory network with self-attention is initialized from an earlier
189-feature model and fine-tuned under a validation-macro-F1 checkpoint rule. On a frozen,
signer-independent internal holdout, motion fusion improves top-1 accuracy from 71.13\% to 71.69\%
and macro F1 from 71.04\% to 71.60\%; mirrored-view macro F1 rises from 69.88\% to 70.98\%.
Top-5 accuracy decreases by 0.34 percentage points, a trade-off reported rather than hidden.
At a validation-selected confidence threshold of 0.60, the system covers 77.50\% of holdout
examples with 84.03\% accuracy among accepted predictions. The user interface adds landmark
quality checks and a stricter provisional operational threshold of 0.71. SignLearn is an isolated
sign classification study and practice aid, not an ASL translator or accessibility service.
\end{abstract}

\noindent\textbf{Keywords:} isolated sign recognition; ASL; landmark sequences; LSTM;
self-attention; selective prediction; human-centered machine learning

\section{Introduction}
Sign languages are complete natural languages, and automatic systems must be described with
linguistic and community-aware care \cite{bragg2019interdisciplinary}. Isolated sign language
recognition (ISLR) maps a bounded video segment to one vocabulary item. It is substantially
narrower than continuous recognition or translation, because the temporal boundary and candidate
vocabulary are given in advance. Consequently, a high score on an isolated benchmark cannot be
interpreted as ASL fluency, translation ability, or real-world accessibility impact.

SignLearn was developed after a deep-learning curriculum to demonstrate the complete path from
data design to a testable application: signer-disjoint splitting, deterministic preprocessing,
temporal modeling, checkpoint selection, comparison against an existing production candidate,
selective prediction, tests, documentation, and a deployable interface. The project asks:
\begin{enumerate}[label=RQ\arabic*.,leftmargin=*]
  \item Can explicit motion and geometry improve a landmark-based temporal baseline?
  \item Can transfer initialization preserve a validated representation while extending its input?
  \item What accuracy--coverage behavior results when the classifier may abstain?
  \item Which claims remain unsupported without a consented external webcam study?
\end{enumerate}

The main contributions are a missingness-aware 492-feature temporal representation; a verified
transfer procedure from a 189-feature baseline; a frozen comparison that reports positive and
negative metric changes; and a quality-gated, confidence-aware deployment prototype with a public
reproducibility package.

\section{Related Work}
\subsection{Isolated sign recognition}
WLASL established a large-scale word-level benchmark and emphasized the difficulty of signer and
viewpoint variation \cite{li2020wlasl}. ASL Citizen later collected community-sourced isolated ASL
videos with attention to realistic participants and use \cite{desai2023aslcitizen}. SignLearn does
not train on or tune against WLASL. It uses a 50-label subset of the Google--Kaggle isolated-sign
competition data \cite{kaggle2023aslsigns}; this narrower scope should remain explicit.

\subsection{Landmark perception and temporal models}
MediaPipe Hands showed that real-time on-device hand tracking can provide structured landmarks
instead of raw pixels \cite{zhang2020mediapipehands}. MediaPipe Holistic jointly exposes face,
pose, and hand landmarks \cite{grishchenko2020holistic,mediapipe2026holisticapi}. This structure
reduces the dimensionality of the learning problem but inherits detector failures and does not
eliminate domain shift.

Long short-term memory networks model sequential dependencies \cite{hochreiter1997lstm}, while
bidirectional wrappers summarize both temporal directions when the entire isolated clip is already
available \cite{tensorflow2024bilstm}. Self-attention provides a learned mechanism for weighting
relationships across time \cite{vaswani2017attention}. SignLearn combines these components because
the task is clip-level classification, not causal streaming recognition.

\subsection{Selective prediction}
Selective classification allows a model to abstain when confidence is insufficient, creating a
risk--coverage trade-off \cite{geifman2017selective}. SelectiveNet integrates a reject option into
training \cite{geifman2019selectivenet}; SignLearn instead applies a post-hoc maximum-probability
threshold. This simpler method does not prove calibration, but it is preferable to presenting every
low-confidence output as reliable.

\section{Data and Evaluation Design}
\subsection{Task and partitions}
The task is single-label classification over 50 isolated signs. Videos are converted to landmark
sequences before model training. Partitions are signer-disjoint so that identity-specific motion
patterns do not occur in both training and evaluation. Development uses training and validation
partitions; a fixed internal holdout is accessed only by the final comparison notebook. Earlier
development inspected a separate Google test resource, so it is not claimed as a pristine external
test. WLASL was not used for training, tuning, checkpoint selection, or promotion.

\subsection{Metrics}
Top-1 accuracy is intuitive but can obscure minority-class failures. Macro F1 gives equal weight to
each label, while top-5 accuracy evaluates candidate ranking. Mirrored-view macro F1 probes camera
orientation robustness. For threshold $\tau$, coverage and selective accuracy are
\begin{align}
\mathrm{coverage}(\tau) &= \frac{1}{N}\sum_{i=1}^{N}
  \mathbb{1}\!\left[\max_c p_{ic}\geq\tau\right],\\
\mathrm{acc}_{\mathrm{accepted}}(\tau) &=
  \frac{\sum_i \mathbb{1}[\hat y_i=y_i]\mathbb{1}[\max_c p_{ic}\geq\tau]}
       {\sum_i \mathbb{1}[\max_c p_{ic}\geq\tau]}.
\end{align}
The implementation follows standard metric definitions \cite{scikit2026metrics}.

\section{Method}
\subsection{Landmarks and normalization}
Each clip is resampled to $T=48$ frames. The base vector $b_t\in\mathbb{R}^{189}$ contains two
hands, selected pose and lip coordinates, and detection indicators. Coordinates are translated and
scaled using shoulder geometry so that body location and camera distance contribute less nuisance
variation. Missing detections are represented through explicit masks; downstream motion and
geometry terms are zeroed when their required landmarks are unavailable.

\subsection{Motion fusion}
For 144 selected coordinate channels, first- and second-order differences are
\begin{align}
v_t &= m_t\odot m_{t-1}\odot(x_t-x_{t-1}),\\
a_t &= m_t\odot m_{t-1}\odot m_{t-2}\odot(v_t-v_{t-1}),
\end{align}
where $m_t$ is a presence mask and $\odot$ denotes element-wise multiplication. A 15-dimensional
geometry vector captures distances involving wrists, shoulders, nose, mouth, hand spread, and
selected fingertips. The complete frame vector is
\[
z_t=[b_t;v_t;a_t;g_t]\in\mathbb{R}^{189+144+144+15}=\mathbb{R}^{492}.
\]
Masking prevents absent landmarks from generating arbitrary jumps or distances.

\subsection{Network}
The sequence $Z\in\mathbb{R}^{48\times492}$ enters a bidirectional LSTM, normalization and dropout,
a second temporal recurrent block, and self-attention. Temporal pooling produces a clip embedding,
followed by dense layers and a 50-way softmax. Dropout regularization follows the general method of
Srivastava et al. \cite{srivastava2014dropout}; layer normalization follows Ba et al.
\cite{ba2016layernorm}; optimization uses Adam \cite{kingma2015adam}.

\subsection{Transfer initialization}
The earlier 189-feature model provides a validated starting point. Its first input kernel is
expanded from 189 to 492 rows: the original rows are copied exactly and new motion/geometry rows
are zero-initialized. Compatible downstream weights are copied unchanged. Before fine-tuning, an
automated audit verifies zero numerical difference between the old model and expanded model when
new channels are zero. This test separates intended feature extension from accidental regression.

\subsection{Checkpoint and promotion policy}
Checkpoints are ranked by validation macro F1, not validation loss. The candidate is compared with
the current production model under frozen metrics. Promotion requires improvement in primary
metrics without a severe robustness or selective-performance failure. This procedure retained an
earlier model when the first candidate failed the gate and promoted the later candidate only after
the final audit.

\section{Experiments and Results}
\subsection{Development progression}
Mirror analysis first established orientation sensitivity. Motion and geometry ablations then
tested the value of explicit temporal channels. A first transfer candidate did not pass the
promotion gate and was rejected. A later run selected by macro F1 passed the frozen audit. These
negative results and decision artifacts remain in the repository.

\subsection{Validation behavior}
The selected candidate reached 73.54\% validation accuracy and 73.34\% validation macro F1. At
$\tau=0.60$, validation coverage was 75.67\% and accuracy among accepted predictions was 85.47\%.
These threshold figures were used for method selection, not as external performance estimates.

\subsection{Frozen internal holdout comparison}
\begin{table}[ht]
\centering
\caption{Previous production model versus the promoted motion-fusion candidate. ``pp'' denotes
percentage points.}
\begin{tabular}{lrrr}
\toprule
Metric & Previous & Motion fusion & Change \\
\midrule
Top-1 accuracy & 71.13\% & \textbf{71.69\%} & +0.56 pp \\
Macro F1 & 71.04\% & \textbf{71.60\%} & +0.56 pp \\
Mirrored macro F1 & 69.88\% & \textbf{70.98\%} & +1.10 pp \\
Top-5 accuracy & \textbf{90.03\%} & 89.69\% & -0.34 pp \\
Accepted accuracy & 83.68\% & \textbf{84.03\%} & +0.35 pp \\
Coverage & 77.43\% & \textbf{77.50\%} & +0.07 pp \\
\bottomrule
\end{tabular}
\end{table}

Motion fusion modestly improves the primary top-1 and macro-F1 metrics and produces a larger gain
under mirrored input. It does not dominate every metric: top-5 accuracy falls by 0.34 points. The
selective estimates use each model's validation-selected threshold; for the promoted model,
$\tau=0.60$. The interface uses a provisional stricter threshold of 0.71 until a separate,
consented webcam calibration study is available.

\subsection{Deployment check}
The repository does not redistribute evaluation recordings. Reviewers can test the application
with a recording they own or are authorized to use. This is an engineering check, not an accuracy
estimate. The Gradio interface supports recorded or uploaded video \cite{gradio2026video}, exposes
top-three candidates, and returns retry guidance when quality or confidence gates fail.

\section{Discussion}
RQ1 is supported within the internal evaluation: adding motion and geometry produces a small
macro-F1 increase and a larger mirrored-view gain. The scale of the improvement argues for
measured language, not a claim that motion fusion solves generalization. RQ2 is supported by the
zero-difference initialization audit and successful fine-tuning. RQ3 shows the expected
accuracy--coverage trade-off: abstention raises accuracy among accepted cases while leaving some
clips unanswered. RQ4 remains open: no claims about real-world user benefit, subgroup parity,
fluency assessment, or accessibility should be made without an externally designed human study.

The explicit representation offers engineering advantages. It is smaller than raw video, easier to
inspect, and capable of masking missing detections. It also discards texture and fine finger detail,
depends on the landmark detector, and may encode signing-space geometry imperfectly. A
bidirectional model is appropriate for already segmented clips but is not a causal continuous-sign
recognizer.

\section{Responsible Use, Limitations, and Threats to Validity}
\paragraph{Construct validity.} The target is a dataset label, not semantic comprehension or ASL
competence. Many signs depend on non-manual markers, context, and continuous coarticulation that a
50-class isolated task does not model.

\paragraph{Internal validity.} Signer-disjoint partitioning reduces identity leakage but does not
remove correlations involving environment, camera, demographics, or collection procedure.
Repeated experimentation on validation data can still overfit model-development choices.

\paragraph{External validity.} No separate consented webcam cohort has been evaluated. The Google
resource was inspected earlier, and WLASL was deliberately reserved and unused. Internal holdout
metrics therefore should not be presented as evidence of broad deployment reliability.

\paragraph{Human and societal risk.} Incorrect predictions may discourage learners or be mistaken
for linguistic assessment. The application must not be used for interpreting, emergencies,
education grading, medical decisions, employment decisions, or access control. Recording another
person requires informed consent. Product decisions and evaluation protocols should involve Deaf
and ASL stakeholders.

\paragraph{Licensing.} The repository's source code is MIT licensed, but that does not automatically
grant redistribution rights for competition data, trained weights, task assets, sample recordings,
or third-party packages. Publishers must complete the model-license checklist before making a
public release.

\section{Future Work}
The highest-priority next step is a preregistered, consented external study with signers and camera
conditions absent from development. Other directions include per-class and subgroup error analysis;
calibration with expected calibration error and reliability plots; landmark-uncertainty features;
causal temporal models for live feedback; vocabulary expansion without collapsing minority-class
performance; and participatory design with Deaf researchers, instructors, and learners.

\section{Conclusion}
SignLearn demonstrates a transparent end-to-end ISLR prototype rather than a translation system.
Transfer-initialized motion fusion raises internal top-1 accuracy and macro F1 by 0.56 percentage
points and mirrored macro F1 by 1.10 points, while slightly reducing top-5 accuracy. Quality gating,
abstention, reproducible artifacts, and explicit limitations make the result more useful as a
portfolio and research project. The central remaining requirement is external human evaluation
designed with the community the technology concerns.

\section*{Artifact Availability}
The accompanying repository contains the inference application, frozen metadata, tests, sample
clips, checksums, decision artifacts, research notebooks, data/model/ethics cards, and both Markdown
and LaTeX versions of this manuscript. Competition data are not redistributed.

\bibliographystyle{plain}
\bibliography{references}
\end{document}
